\paragraph{Perturbation--sensitivity dissociation and cascade dissolution.}
We hypothesize that the dissociation reflects LLM-specific cascade dynamics: each acknowledgment alters shared history via autoregressive conditioning, raising the next agent's ack probability (Qwen 90.7\%, Llama 61.9\%, both $p < 10^{-13}$; Appendix~\ref{app:cascade}). Aggregate removal dissolves these chains rather than accumulating errors, consistent with flat dose-response curves (Appendix~\ref{app:dose_response}). Position-controlled logistic regression confirms genuine sequential dependence (previous-ack OR $= 6.43$ Qwen, $2.02$ Llama, both $p < 10^{-5}$; Appendix~\ref{app:cascade_position_control}). Notably, the position effect itself reverses across backbones (Qwen OR $= 1.95$, amplifying; Llama OR $= 0.73$, dampening), yet sequential dependence remains significant in both---ruling out a pure positional confound. The mechanism is ack-specific (FG degrades accuracy, $p = 0.031$, exploratory) and distinct from human backchannels~\cite{bavelas_co_narrators_2000}; the double dissociation (\S\ref{sec:impact_gradient}) argues against a noise hypothesis.
We conjecture that aggregate safety holds when: (i)~cascade rate $P(\text{ack} \mid \text{ack}) > 0.5$ and (ii)~bigram novelty below 10\%. Qwen satisfies both (90.7\%, 5.7\%); Llama satisfies~(i) but not~(ii), with safety attributed to error-channel removal (Appendix~\ref{app:llama_mechanism}). Both conditions are violated in MATH where pruning degrades accuracy. Causal verification through controlled intervention remains needed. These thresholds are post-hoc observations derived from two backbone families; prospective validation on additional architectures is needed before they can serve as deployment guidelines.

\paragraph{Dual-mechanism robustness.}
The pruning safety result is robust across two mechanistically distinct pathways. For Qwen, where acknowledgments are terse (median 4 tokens, 5.7\% bigram novelty), pruning removes redundant structural signals---a \emph{redundancy removal} mechanism. For Llama, where acknowledgments are verbose (mean 59.8 tokens, 65.6\% bigram novelty) yet carry zero corrections and non-zero error injection (Fisher $p = 0.004$), pruning eliminates a one-directional \emph{error channel}. That the dissociation holds across both mechanisms---arising from redundancy removal in one backbone and error-channel removal in the other---suggests a structural property of acknowledgment cascades rather than a mechanism-specific artifact. For Llama on GSM8K, the heuristic CI is entirely positive ($[+0.41, +3.99]$\,pp, 5~seeds)---consistent with error-channel removal improving downstream signal quality---indicating that pruning is not merely safe but actively beneficial; on StrategyQA, heuristic pruning remains non-inferior (CI $[-0.49, +6.11]$\,pp, 8~seeds) with higher seed-level variance, suggesting the magnitude of the effect is domain-modulated. Temperature robustness mirrors the typology distinction: Qwen's function-organized cascade persists under stochastic decoding (92.2\% at $t{=}0.7$ vs.\ 90.7\% greedy), while Llama's undifferentiated pattern is temperature-sensitive (suppress drops from 37\% to 28\%). This dual-mechanism evidence strengthens the case for the analysis protocol as the primary contribution: the protocol detects prunability regardless of the underlying mechanism.

\paragraph{Prompt independence.}
Removing the explicit confirmation instruction from the checker prompt leaves the heuristic ack detection rate at 76.2\% (Qwen/GSM8K, $N{=}100$, char$<$30 heuristic), compared to 47.9\% oracle ack rate under the standard prompt. The 76.2\% figure reflects heuristic sensitivity to short messages, not oracle ack prevalence; the two are not directly comparable (different detection methods, different prompt conditions). Regardless, ack-like behavior persists without prompting, suggesting acknowledgments are emergent rather than prompt-driven. Primary pruning safety evidence (\S\ref{sec:e2e_pruning}) is based entirely on oracle labels, not heuristic detection.

\paragraph{Communication structure and validation.}
Backbone dominates communication typology ($V = 0.45$ vs.\ $0.18$; \S\ref{sec:cross_domain}), with cross-model concordance confirming model-specific ack placement (stratified $\kappa = 0.018$; Appendix~\ref{app:cross_model_concordance}). Suppression rates (36--56\%) translate to equivalent inference-call reductions (Appendix~\ref{app:efficiency}). Four-annotator validation confirms taxonomy reliability (pairwise $\kappa = 0.76$, 5-class; binary ack/non-ack $\kappa = 0.81$; Appendix~\ref{app:kappa_tiers}); Gemma's $\kappa = 0.35$ falls below threshold and its results are treated as exploratory. Acknowledgment prevalence varies sharply with topology: debate produces near-zero acks (0.2\%), and code-generation (HumanEval) only 4\%, suggesting acknowledgments are an emergent property of solver-checker interaction rather than a universal feature of multi-agent communication.

\paragraph{Detector limitations.}
Heuristic detectors are proof-of-concept classifiers with non-transferability across domains (F1: $0.845 \to 0.213$), backbones (Gemma $\kappa = 0.35$), and generations (Qwen3-8B: $-14$\,pp; \S\ref{sec:boundary}). No universal length proxy exists ($56\times$ coverage gap). Oracle pruning remains robust across all tested cells, confirming the detector---not the pruning principle---is the limiting factor. An adversary who can inject acknowledgment-patterned messages into a multi-agent pipeline could exploit the cascade dynamics to cause outsized behavioral disruption with minimal detectable content; defenses should monitor structural patterns alongside content for integrity.

\paragraph{Limitations.}
Oracle labels face a distribution-shift caveat (offline vs.\ online gap: negligible for Qwen, larger for Llama; Appendix~\ref{app:online}). Pruning safety is backbone- and domain-specific (\S\ref{sec:boundary}); all experiments use two-agent topologies under greedy decoding; scale evidence at 72B is preliminary (single seed). Additionally, annotation content previews for StrategyQA, ARC, and MATH are predominantly short ($< 100$ characters), so ack/non-ack binary labels---where acknowledgments are inherently brief---are unlikely affected, but fine-grained 5-class annotations on longer messages may suffer information loss from truncation. The learned DistilBERT classifier (F1~$= 0.867$; \S\ref{sec:learned_classifier}) has not been tested in the online pruning loop; heuristic detectors degrade across domains (F1: $0.845 \to 0.213$). Bridging oracle pruning safety to practical deployment---via online evaluation of learned classifiers---is immediate future work. We recommend empirical screening ($N \geq 200$, ${\geq}\,80\%$ power for ${\geq}\,5$\,pp effects) before deployment; extended limitations appear in Appendix~\ref{app:extended_limitations}.
